\documentclass[pdflatex,sn-mathphys-num]{sn-jnl}

\usepackage{graphicx}%
\usepackage{multirow}%
\usepackage{amsmath,amssymb,amsfonts}%
\usepackage{amsthm}%
\usepackage{mathrsfs}%
\usepackage[title]{appendix}%
\usepackage{xcolor}%
\usepackage{textcomp}%
\usepackage{manyfoot}%
\usepackage{booktabs}%
\usepackage{algorithm}%
\usepackage{algorithmicx}%
\usepackage{algpseudocode}%
\usepackage{listings}%
\usepackage{lmodern}%
\usepackage{siunitx}
\usepackage{microtype}
\usepackage{derivative}
\usepackage[version=4]{mhchem}
\usepackage{cleveref}
\usepackage[spanish]{babel}

\renewcommand{\arraystretch}{1.2}
\sisetup{group-digits=true,
  group-separator={\,},
  separate-uncertainty
}

\theoremstyle{thmstyleone}%
\newtheorem{theorem}{Teorema}%
\newtheorem{proposition}[theorem]{Proposición}%

\theoremstyle{thmstyletwo}%
\newtheorem{example}{Ejemplo}%
\newtheorem{remark}{Observación}%

\theoremstyle{thmstylethree}%
\newtheorem{definition}{Definición}%

\DeclareSIUnit\angstrom{\text{\r{A}}}

\raggedbottom
\unnumbered% uncomment this for unnumbered level heads

\begin{document}

\title[Diseño de Nanomateriales]{Diseño de Nanomateriales y Teoría del Funcional
de la Densidad: Resumen de la Charla de Sindy J. Rodriguez}

\author{\fnm{Julian L.}
\sur{Avila-Martinez}}\email{jlavilam@udistrital.edu.co}

\affil{\orgdiv{Programa Académico de Física}, \orgname{Universidad Distrital
Francisco José de Caldas}, \orgaddress{\city{Bogotá D.C.},
\country{Colombia}}}

\abstract{
  Se expone la sinergia entre el modelamiento computacional y la caracterización
  experimental en el diseño de nanomateriales. Se analiza la reducción de
  dimensionalidad del problema de muchos cuerpos a través de la Teoría del
  Funcional de la Densidad (DFT) y su correspondencia con observaciones
  empíricas mediante microscopía de efecto túnel (STM). Adicionalmente, se
  examina la dinámica de intercalación para el almacenamiento de energía y la
  funcionalización de superficies para modular el transporte electrónico
  multidimensional.
}

\keywords{Teoría del funcional de la densidad, nanomateriales, microscopía de
efecto túnel, física de superficies, hamiltonianos de transporte}

\maketitle

\section{Reducción de Dimensionalidad y el Hamiltoniano Cuántico}
\label{sec:reduccion}

La resolución directa de la ecuación de Schrödinger independiente del tiempo
para sistemas de muchos cuerpos presenta un desafío computacional intratable
debido a la explosión combinatoria de los grados de libertad.  Para un sistema
compuesto por $N$ partículas, la función de onda $\Psi$ reside en un espacio de
configuración abstracto de $3N$ dimensiones. Como se ilustró en la ponencia, la
discretización espacial exhaustiva de un sistema diatómico simple exigiría
recursos computacionales y tiempos de procesamiento que superan escalas
cosmológicas.

Para sortear esta catástrofe de dimensionalidad, se recurre a la Teoría del
Funcional de la Densidad (DFT). Este marco teórico establece que el estado
fundamental de un sistema cuántico está unívocamente determinado por su densidad
electrónica espacial, $n(\mathbf{r})$. Esta formulación mapea el problema
original a un espacio vectorial de estrictamente tres dimensiones, permitiendo
la minimización del funcional de energía sin requerir el cálculo explícito de la
hiperfunción de onda electrónica. No obstante, la evaluación de las métricas
dinámicas fuera del equilibrio, como el transporte electrónico impulsado por
diferencias de potencial, sigue exigiendo acoplamientos no triviales con
funciones de Green para mitigar el costo del análisis.

\section{Física de Superficies y Mapeo por Efecto Túnel}
\label{sec:superficies}

A escala nanométrica, las propiedades macroscópicas de los cristales están
dictadas por su simetría topológica y la periodicidad espacial de su red
atómica, con independencia estricta de su homogeneidad química. La dramática
divergencia en el comportamiento conductor y en la sección eficaz óptica entre
el diamante, el grafito y el grafeno evidencia cómo la geometría del retículo
gobierna los autovalores energéticos del sistema.

La validación empírica de estas arquitecturas y de los mecanismos de adsorción
se realiza mediante instrumentación en condiciones de ultra alto vacío,
notablemente mediante la microscopía de efecto túnel (STM). Este método no mide
la topografía geométrica clásica, sino que explota el tunelamiento de la función
de onda electrónica a través de una barrera de potencial hacia el sustrato. La
corriente eléctrica mapea la densidad local de estados (LDOS), revelando los
dominios espaciales de alta probabilidad electrónica. El complemento analítico
recae en la espectroscopía de fotoelectrones de rayos X (XPS), indispensable
para resolver el grado de hibridación orbital y la pureza química en la
estructura de bandas.

\section{Dinámica de Intercalación y Modulación del Transporte}
\label{sec:intercalacion}

La funcionalización de superficies de baja dimensionalidad actúa como un
operador que perturba el hamiltoniano del material anfitrión, modificando el
transporte intrínseco mediante dopaje directo o ensamblaje supramolecular.
Análisis teóricos demuestran que la quimisorción de cadenas de aminoácidos sobre
redes bidimensionales de grafeno induce asimetrías severas en las curvas de
corriente-voltaje; dichas perturbaciones locales imponen un comportamiento
análogo al de un transistor molecular de rectificación direccional.

En el contexto termodinámico del almacenamiento de energía, se contrastó
matemáticamente el mecanismo de intercalación del cloruro de aluminio
(\ce{AlCl4-}) frente a los convencionales iones de litio en redes de grafito.
Las predicciones teóricas concluyeron que las moléculas intercalantes de
aluminio franquean menores barreras de difusión y presentan estabilidades
estructurales competitivas frente a sus contrapartes iónicas. La simulación de
interacciones intermoleculares evidenció un comportamiento altamente no lineal
donde la penetración escapa a distribuciones gaussianas isotrópicas; en su
lugar, se agregan bajo mecanismos de etapas mixtas (mixed-stage) creando
distribuciones espacialmente discretas de clústeres. Esta estratificación local
induce tensores de estrés significativos sobre la red cristalina, lo que limita
la invariancia estructural y condiciona la reversibilidad funcional del
material.

\end{document}


\end{document}
